% --- Archon physics formalization source begin ---
% archon:physics
% archon:covers PhyXMiniProblems/problem_phyx_mini_0485.lean
% archon:source-report reports/phyx_mini/problem_phyx_mini_0485.source.json

\chapter{Physics problem phyx\_mini\_0485}
\label{ch:PhyXMiniProblems_problem_phyx_mini_0485}

\paragraph{Problem source.}
\#\# Physical scenario

A \textbackslash{}(0.40\textbackslash{},\textbackslash{}mathrm\{kg\}\textbackslash{}) iron horseshoe, just forged and very hot (figure), is dropped into \textbackslash{}(1.25\textbackslash{},\textbackslash{}mathrm\{L\}\textbackslash{}) of water in a \textbackslash{}(0.30\textbackslash{},\textbackslash{}mathrm\{kg\}\textbackslash{}) iron pot initially at \textbackslash{}(20.0\textasciicircum{}\textbackslash{}circ\textbackslash{}mathrm\{C\}\textbackslash{}).The final equilibrium temperature is \textbackslash{}(25.0\textasciicircum{}\textbackslash{}circ\textbackslash{}mathrm\{C\}\textbackslash{}).

\#\# Question

Estimate the initial temperature of the hot horseshoe.

\#\# Answer choices

- A: A: 140°C

- B: B: 150°C

- C: C: 160°C

- D: D: 170°C

\#\# Figure caption (auxiliary; use the image as primary evidence)

The image features a close-up scene of blacksmithing activity. A glowing, hot horseshoe, colored orange from the heat, is held in place with tongs by a pair of hands. The horseshoe is positioned against an anvil, a sturdy block often used as a surface for forging and shaping metal. The background is softly blurred, keeping the focus on the horseshoe and anvil interaction. There is a sense of movement and heat, but no text is present in the image.

\#\# Dataset metadata

Specific Heat and Calorimetry; Multi-Formula Reasoning, Physical Model Grounding Reasoning

\paragraph{Recorded answer/context.}
D: D: 170°C

\paragraph{Figure/image path.}
/root/proposal\_for\_physic/science-mango/phyx\_mini\_run/phyx\_data/test\_image/485.png

\paragraph{Formalization target.}
create a compiling Lean file with sorry bodies at `PhyXMiniProblems/problem\_phyx\_mini\_0485.lean`. The Lean declarations must preserve the physical quantities, dimensions or dimensional roles, figure labels, governing-law hypotheses, and final relation expressed by this problem.
Use Mathlib/Physlib names found through LeanExplore where available. If a physics API is missing, introduce faithful local abstractions rather than scalar placeholder aliases.

\begin{theorem}[Physics formalization target]
\label{thm:physics:phyx_mini_0485:target}
\lean{PhyXMiniProblems.ProblemPhyXMini0485.problem_phyx_mini_0485}
\uses{def:physics:phyx-mini-0485:phyxminiproblems-problemphyxmini0485-horseshoewatercalorimetrysetup, def:physics:phyx-mini-0485:phyxminiproblems-problemphyxmini0485-initialhorseshoetemperatureindegreescelsius, def:physics:phyx-mini-0485:phyxminiproblems-problemphyxmini0485-matcheshorseshoewaterscenario, def:physics:phyx-mini-0485:phyxminiproblems-problemphyxmini0485-matchesproblemreadouts, def:physics:phyx-mini-0485:phyxminiproblems-problemphyxmini0485-matchessuppliedblacksmithfigure, def:physics:phyx-mini-0485:phyxminiproblems-problemphyxmini0485-usesreferenceironandwaterthermaldata, def:physics:phyx-mini-0485:phyxminiproblems-problemphyxmini0485-hasphysicalhorseshoewaterparameters, def:physics:phyx-mini-0485:phyxminiproblems-problemphyxmini0485-satisfieswaterbulkmasslaw, def:physics:phyx-mini-0485:phyxminiproblems-problemphyxmini0485-satisfiesisolatedhorseshoewatercalorimetry, def:physics:phyx-mini-0485:phyxminiproblems-problemphyxmini0485-recordeddatasetanswer, def:physics:phyx-mini-0485:phyxminiproblems-problemphyxmini0485-roundstodisplayedtendegrees, def:physics:phyx-mini-0485:phyxminiproblems-problemphyxmini0485-isuniqueclosestdisplayedtemperature}
The assigned autoformalize agent should translate this physics problem into Lean declarations in the covered file, with theorem and lemma proof bodies written as `by sorry`.
\end{theorem}
\begin{proof}
This is an autoformalization task, not a proof task. Produce statements that can later be proved by the physics prover without weakening the physical model.
\end{proof}

% --- Archon declaration topology begin ---
\section{Declaration topology}
\begin{definition}[volume Dimension]
  \label{def:physics:phyx-mini-0485:phyxminiproblems-problemphyxmini0485-volumedimension}
  \lean{PhyXMiniProblems.ProblemPhyXMini0485.volumeDimension}
  The physical dimension of volume, L³
\end{definition}
\begin{proof}
  This is the declared model definition of volume Dimension.
\end{proof}
\begin{definition}[mass Density Dimension]
  \label{def:physics:phyx-mini-0485:phyxminiproblems-problemphyxmini0485-massdensitydimension}
  \lean{PhyXMiniProblems.ProblemPhyXMini0485.massDensityDimension}
  The physical dimension of mass density, M L⁻³
\end{definition}
\begin{proof}
  This is the declared model definition of mass Density Dimension.
\end{proof}
\begin{definition}[specific Heat Capacity Dimension]
  \label{def:physics:phyx-mini-0485:phyxminiproblems-problemphyxmini0485-specificheatcapacitydimension}
  \lean{PhyXMiniProblems.ProblemPhyXMini0485.specificHeatCapacityDimension}
  The physical dimension of mass-specific heat capacity, L² T⁻² Θ⁻¹
\end{definition}
\begin{proof}
  This is the declared model definition of specific Heat Capacity Dimension.
\end{proof}
\begin{definition}[Mass Quantity]
  \label{def:physics:phyx-mini-0485:phyxminiproblems-problemphyxmini0485-massquantity}
  \lean{PhyXMiniProblems.ProblemPhyXMini0485.MassQuantity}
  A nonnegative physical mass, independent of a choice of units
\end{definition}
\begin{proof}
  This is the declared model definition of Mass Quantity.
\end{proof}
\begin{definition}[Volume Quantity]
  \label{def:physics:phyx-mini-0485:phyxminiproblems-problemphyxmini0485-volumequantity}
  \lean{PhyXMiniProblems.ProblemPhyXMini0485.VolumeQuantity}
  \uses{def:physics:phyx-mini-0485:phyxminiproblems-problemphyxmini0485-volumedimension}
  A nonnegative physical volume, independent of a choice of units
\end{definition}
\begin{proof}
  This is the declared model definition of Volume Quantity.
\end{proof}
\begin{definition}[Mass Density Quantity]
  \label{def:physics:phyx-mini-0485:phyxminiproblems-problemphyxmini0485-massdensityquantity}
  \lean{PhyXMiniProblems.ProblemPhyXMini0485.MassDensityQuantity}
  \uses{def:physics:phyx-mini-0485:phyxminiproblems-problemphyxmini0485-massdensitydimension}
  A nonnegative physical mass density, independent of a choice of units
\end{definition}
\begin{proof}
  This is the declared model definition of Mass Density Quantity.
\end{proof}
\begin{definition}[Specific Heat Capacity Quantity]
  \label{def:physics:phyx-mini-0485:phyxminiproblems-problemphyxmini0485-specificheatcapacityquantity}
  \lean{PhyXMiniProblems.ProblemPhyXMini0485.SpecificHeatCapacityQuantity}
  \uses{def:physics:phyx-mini-0485:phyxminiproblems-problemphyxmini0485-specificheatcapacitydimension}
  A nonnegative mass-specific heat capacity, independent of a choice of units
\end{definition}
\begin{proof}
  This is the declared model definition of Specific Heat Capacity Quantity.
\end{proof}
\begin{definition}[mass Readout]
  \label{def:physics:phyx-mini-0485:phyxminiproblems-problemphyxmini0485-massreadout}
  \lean{PhyXMiniProblems.ProblemPhyXMini0485.massReadout}
  \uses{def:physics:phyx-mini-0485:phyxminiproblems-problemphyxmini0485-massquantity}
  Read a physical mass in a selected mass unit
\end{definition}
\begin{proof}
  This is the declared model definition of mass Readout.
\end{proof}
\begin{definition}[volume Readout]
  \label{def:physics:phyx-mini-0485:phyxminiproblems-problemphyxmini0485-volumereadout}
  \lean{PhyXMiniProblems.ProblemPhyXMini0485.volumeReadout}
  \uses{def:physics:phyx-mini-0485:phyxminiproblems-problemphyxmini0485-volumequantity}
  Read a physical volume in the cube of a selected length unit
\end{definition}
\begin{proof}
  This is the declared model definition of volume Readout.
\end{proof}
\begin{definition}[mass Density Readout]
  \label{def:physics:phyx-mini-0485:phyxminiproblems-problemphyxmini0485-massdensityreadout}
  \lean{PhyXMiniProblems.ProblemPhyXMini0485.massDensityReadout}
  \uses{def:physics:phyx-mini-0485:phyxminiproblems-problemphyxmini0485-massdensityquantity}
  Read density in selected mass-per-cubic-length units
\end{definition}
\begin{proof}
  This is the declared model definition of mass Density Readout.
\end{proof}
\begin{definition}[mass In Kilograms]
  \label{def:physics:phyx-mini-0485:phyxminiproblems-problemphyxmini0485-massinkilograms}
  \lean{PhyXMiniProblems.ProblemPhyXMini0485.massInKilograms}
  \uses{def:physics:phyx-mini-0485:phyxminiproblems-problemphyxmini0485-massquantity, def:physics:phyx-mini-0485:phyxminiproblems-problemphyxmini0485-massreadout}
  Kilogram readout used by the calorimetry law
\end{definition}
\begin{proof}
  This is the declared model definition of mass In Kilograms.
\end{proof}
\begin{definition}[volume In Cubic Meters]
  \label{def:physics:phyx-mini-0485:phyxminiproblems-problemphyxmini0485-volumeincubicmeters}
  \lean{PhyXMiniProblems.ProblemPhyXMini0485.volumeInCubicMeters}
  \uses{def:physics:phyx-mini-0485:phyxminiproblems-problemphyxmini0485-volumequantity, def:physics:phyx-mini-0485:phyxminiproblems-problemphyxmini0485-volumereadout}
  Cubic-metre readout of a physical volume
\end{definition}
\begin{proof}
  This is the declared model definition of volume In Cubic Meters.
\end{proof}
\begin{definition}[volume In Liters]
  \label{def:physics:phyx-mini-0485:phyxminiproblems-problemphyxmini0485-volumeinliters}
  \lean{PhyXMiniProblems.ProblemPhyXMini0485.volumeInLiters}
  \uses{def:physics:phyx-mini-0485:phyxminiproblems-problemphyxmini0485-volumequantity, def:physics:phyx-mini-0485:phyxminiproblems-problemphyxmini0485-volumeincubicmeters}
  Litre readout, using 1 m³ = 1000 L
\end{definition}
\begin{proof}
  This is the declared model definition of volume In Liters.
\end{proof}
\begin{definition}[mass Density In Kilograms Per Cubic Meter]
  \label{def:physics:phyx-mini-0485:phyxminiproblems-problemphyxmini0485-massdensityinkilogramspercubicmeter}
  \lean{PhyXMiniProblems.ProblemPhyXMini0485.massDensityInKilogramsPerCubicMeter}
  \uses{def:physics:phyx-mini-0485:phyxminiproblems-problemphyxmini0485-massdensityquantity, def:physics:phyx-mini-0485:phyxminiproblems-problemphyxmini0485-massdensityreadout}
  SI density readout in kilograms per cubic metre
\end{definition}
\begin{proof}
  This is the declared model definition of mass Density In Kilograms Per Cubic Meter.
\end{proof}
\begin{definition}[mass Density In Kilograms Per Liter]
  \label{def:physics:phyx-mini-0485:phyxminiproblems-problemphyxmini0485-massdensityinkilogramsperliter}
  \lean{PhyXMiniProblems.ProblemPhyXMini0485.massDensityInKilogramsPerLiter}
  \uses{def:physics:phyx-mini-0485:phyxminiproblems-problemphyxmini0485-massdensityquantity, def:physics:phyx-mini-0485:phyxminiproblems-problemphyxmini0485-massdensityinkilogramspercubicmeter}
  Kilogram-per-litre density readout, using 1 m³ = 1000 L
\end{definition}
\begin{proof}
  This is the declared model definition of mass Density In Kilograms Per Liter.
\end{proof}
\begin{definition}[specific Heat In Joules Per Kilogram Kelvin]
  \label{def:physics:phyx-mini-0485:phyxminiproblems-problemphyxmini0485-specificheatinjoulesperkilogramkelvin}
  \lean{PhyXMiniProblems.ProblemPhyXMini0485.specificHeatInJoulesPerKilogramKelvin}
  \uses{def:physics:phyx-mini-0485:phyxminiproblems-problemphyxmini0485-specificheatcapacityquantity}
  SI readout of mass-specific heat capacity in joules per kilogram-kelvin
\end{definition}
\begin{proof}
  This is the declared model definition of specific Heat In Joules Per Kilogram Kelvin.
\end{proof}
\begin{definition}[heat In Joules]
  \label{def:physics:phyx-mini-0485:phyxminiproblems-problemphyxmini0485-heatinjoules}
  \lean{PhyXMiniProblems.ProblemPhyXMini0485.heatInJoules}
  SI readout of a signed heat transfer in joules
\end{definition}
\begin{proof}
  This is the declared model definition of heat In Joules.
\end{proof}
\begin{definition}[Celsius Temperature Reading]
  \label{def:physics:phyx-mini-0485:phyxminiproblems-problemphyxmini0485-celsiustemperaturereading}
  \lean{PhyXMiniProblems.ProblemPhyXMini0485.CelsiusTemperatureReading}
  An affine Celsius coordinate attached to a Physlib absolute temperature. The offset is 273.15 K = 5463/20 K
\end{definition}
\begin{proof}
  This is the declared model definition of Celsius Temperature Reading.
\end{proof}
\begin{definition}[Thermal Body]
  \label{def:physics:phyx-mini-0485:phyxminiproblems-problemphyxmini0485-thermalbody}
  \lean{PhyXMiniProblems.ProblemPhyXMini0485.ThermalBody}
  The three bodies that exchange heat
\end{definition}
\begin{proof}
  This is the declared model definition of Thermal Body.
\end{proof}
\begin{definition}[Thermal Material]
  \label{def:physics:phyx-mini-0485:phyxminiproblems-problemphyxmini0485-thermalmaterial}
  \lean{PhyXMiniProblems.ProblemPhyXMini0485.ThermalMaterial}
  The two thermal materials named in the problem
\end{definition}
\begin{proof}
  This is the declared model definition of Thermal Material.
\end{proof}
\begin{definition}[Exterior Boundary Condition]
  \label{def:physics:phyx-mini-0485:phyxminiproblems-problemphyxmini0485-exteriorboundarycondition}
  \lean{PhyXMiniProblems.ProblemPhyXMini0485.ExteriorBoundaryCondition}
  Thermal condition at the boundary of the modeled three-body system
\end{definition}
\begin{proof}
  This is the declared model definition of Exterior Boundary Condition.
\end{proof}
\begin{definition}[Horseshoe Condition]
  \label{def:physics:phyx-mini-0485:phyxminiproblems-problemphyxmini0485-horseshoecondition}
  \lean{PhyXMiniProblems.ProblemPhyXMini0485.HorseshoeCondition}
  Manufacturing and thermal condition of the horseshoe before immersion
\end{definition}
\begin{proof}
  This is the declared model definition of Horseshoe Condition.
\end{proof}
\begin{definition}[Horseshoe Water Contact]
  \label{def:physics:phyx-mini-0485:phyxminiproblems-problemphyxmini0485-horseshoewatercontact}
  \lean{PhyXMiniProblems.ProblemPhyXMini0485.HorseshoeWaterContact}
  Contact relation between the horseshoe and water
\end{definition}
\begin{proof}
  This is the declared model definition of Horseshoe Water Contact.
\end{proof}
\begin{definition}[Figure Object]
  \label{def:physics:phyx-mini-0485:phyxminiproblems-problemphyxmini0485-figureobject}
  \lean{PhyXMiniProblems.ProblemPhyXMini0485.FigureObject}
  Objects identified in the supplied blacksmithing image and caption
\end{definition}
\begin{proof}
  This is the declared model definition of Figure Object.
\end{proof}
\begin{definition}[Horseshoe Glow Color]
  \label{def:physics:phyx-mini-0485:phyxminiproblems-problemphyxmini0485-horseshoeglowcolor}
  \lean{PhyXMiniProblems.ProblemPhyXMini0485.HorseshoeGlowColor}
  Recorded visible glow color of the hot horseshoe
\end{definition}
\begin{proof}
  This is the declared model definition of Horseshoe Glow Color.
\end{proof}
\begin{definition}[Supplied Blacksmith Figure]
  \label{def:physics:phyx-mini-0485:phyxminiproblems-problemphyxmini0485-suppliedblacksmithfigure}
  \lean{PhyXMiniProblems.ProblemPhyXMini0485.SuppliedBlacksmithFigure}
  \uses{def:physics:phyx-mini-0485:phyxminiproblems-problemphyxmini0485-figureobject, def:physics:phyx-mini-0485:phyxminiproblems-problemphyxmini0485-horseshoeglowcolor}
  Qualitative evidence from the supplied image. It contains no printed temperature, mass, volume, material constant, or answer label
\end{definition}
\begin{proof}
  This is the declared model definition of Supplied Blacksmith Figure.
\end{proof}
\begin{definition}[Horseshoe Water Calorimetry Setup]
  \label{def:physics:phyx-mini-0485:phyxminiproblems-problemphyxmini0485-horseshoewatercalorimetrysetup}
  \lean{PhyXMiniProblems.ProblemPhyXMini0485.HorseshoeWaterCalorimetrySetup}
  \uses{def:physics:phyx-mini-0485:phyxminiproblems-problemphyxmini0485-massquantity, def:physics:phyx-mini-0485:phyxminiproblems-problemphyxmini0485-volumequantity, def:physics:phyx-mini-0485:phyxminiproblems-problemphyxmini0485-massdensityquantity, def:physics:phyx-mini-0485:phyxminiproblems-problemphyxmini0485-specificheatcapacityquantity, def:physics:phyx-mini-0485:phyxminiproblems-problemphyxmini0485-celsiustemperaturereading, def:physics:phyx-mini-0485:phyxminiproblems-problemphyxmini0485-thermalbody, def:physics:phyx-mini-0485:phyxminiproblems-problemphyxmini0485-thermalmaterial, def:physics:phyx-mini-0485:phyxminiproblems-problemphyxmini0485-exteriorboundarycondition, def:physics:phyx-mini-0485:phyxminiproblems-problemphyxmini0485-horseshoecondition, def:physics:phyx-mini-0485:phyxminiproblems-problemphyxmini0485-horseshoewatercontact, def:physics:phyx-mini-0485:phyxminiproblems-problemphyxmini0485-suppliedblacksmithfigure}
  Independent physical quantities and qualitative roles of the experiment. Signed body heats are positive into each body; the surroundings heat is positive into the modeled three-body system
\end{definition}
\begin{proof}
  This is the declared model definition of Horseshoe Water Calorimetry Setup.
\end{proof}
\begin{definition}[initial Horseshoe Temperature In Degrees Celsius]
  \label{def:physics:phyx-mini-0485:phyxminiproblems-problemphyxmini0485-initialhorseshoetemperatureindegreescelsius}
  \lean{PhyXMiniProblems.ProblemPhyXMini0485.initialHorseshoeTemperatureInDegreesCelsius}
  \uses{def:physics:phyx-mini-0485:phyxminiproblems-problemphyxmini0485-horseshoewatercalorimetrysetup}
  The requested Celsius coordinate, exposed without assigning it a value
\end{definition}
\begin{proof}
  This is the declared model definition of initial Horseshoe Temperature In Degrees Celsius.
\end{proof}
\begin{definition}[Matches Horseshoe Water Scenario]
  \label{def:physics:phyx-mini-0485:phyxminiproblems-problemphyxmini0485-matcheshorseshoewaterscenario}
  \lean{PhyXMiniProblems.ProblemPhyXMini0485.MatchesHorseshoeWaterScenario}
  \uses{def:physics:phyx-mini-0485:phyxminiproblems-problemphyxmini0485-horseshoewatercalorimetrysetup, def:physics:phyx-mini-0485:phyxminiproblems-problemphyxmini0485-initialhorseshoetemperatureindegreescelsius}
  Qualitative process conditions stated or selected by the textbook model
\end{definition}
\begin{proof}
  This is the declared model definition of Matches Horseshoe Water Scenario.
\end{proof}
\begin{definition}[Matches Problem Readouts]
  \label{def:physics:phyx-mini-0485:phyxminiproblems-problemphyxmini0485-matchesproblemreadouts}
  \lean{PhyXMiniProblems.ProblemPhyXMini0485.MatchesProblemReadouts}
  \uses{def:physics:phyx-mini-0485:phyxminiproblems-problemphyxmini0485-massinkilograms, def:physics:phyx-mini-0485:phyxminiproblems-problemphyxmini0485-volumeinliters, def:physics:phyx-mini-0485:phyxminiproblems-problemphyxmini0485-horseshoewatercalorimetrysetup}
  Numerical data printed in the problem. The unknown initial horseshoe temperature and the recorded answer do not occur here
\end{definition}
\begin{proof}
  This is the declared model definition of Matches Problem Readouts.
\end{proof}
\begin{definition}[Matches Supplied Blacksmith Figure]
  \label{def:physics:phyx-mini-0485:phyxminiproblems-problemphyxmini0485-matchessuppliedblacksmithfigure}
  \lean{PhyXMiniProblems.ProblemPhyXMini0485.MatchesSuppliedBlacksmithFigure}
  \uses{def:physics:phyx-mini-0485:phyxminiproblems-problemphyxmini0485-suppliedblacksmithfigure}
  Qualitative facts recorded from the supplied blacksmithing image
\end{definition}
\begin{proof}
  This is the declared model definition of Matches Supplied Blacksmith Figure.
\end{proof}
\begin{definition}[Uses Reference Iron And Water Thermal Data]
  \label{def:physics:phyx-mini-0485:phyxminiproblems-problemphyxmini0485-usesreferenceironandwaterthermaldata}
  \lean{PhyXMiniProblems.ProblemPhyXMini0485.UsesReferenceIronAndWaterThermalData}
  \uses{def:physics:phyx-mini-0485:phyxminiproblems-problemphyxmini0485-massdensityinkilogramsperliter, def:physics:phyx-mini-0485:phyxminiproblems-problemphyxmini0485-specificheatinjoulesperkilogramkelvin, def:physics:phyx-mini-0485:phyxminiproblems-problemphyxmini0485-horseshoewatercalorimetrysetup}
  Rounded textbook reference data needed for the estimate: water density is 1 kg/L, iron specific heat is 450 J/(kg K), and water specific heat is 4186 J/(kg K). These values are not printed in the problem, so they are kept separate from MatchesProblemReadouts
\end{definition}
\begin{proof}
  This is the declared model definition of Uses Reference Iron And Water Thermal Data.
\end{proof}
\begin{definition}[Has Physical Horseshoe Water Parameters]
  \label{def:physics:phyx-mini-0485:phyxminiproblems-problemphyxmini0485-hasphysicalhorseshoewaterparameters}
  \lean{PhyXMiniProblems.ProblemPhyXMini0485.HasPhysicalHorseshoeWaterParameters}
  \uses{def:physics:phyx-mini-0485:phyxminiproblems-problemphyxmini0485-massinkilograms, def:physics:phyx-mini-0485:phyxminiproblems-problemphyxmini0485-volumeinliters, def:physics:phyx-mini-0485:phyxminiproblems-problemphyxmini0485-massdensityinkilogramsperliter, def:physics:phyx-mini-0485:phyxminiproblems-problemphyxmini0485-specificheatinjoulesperkilogramkelvin, def:physics:phyx-mini-0485:phyxminiproblems-problemphyxmini0485-horseshoewatercalorimetrysetup}
  Positivity conditions selecting physically meaningful quantities
\end{definition}
\begin{proof}
  This is the declared model definition of Has Physical Horseshoe Water Parameters.
\end{proof}
\begin{definition}[Satisfies Water Bulk Mass Law]
  \label{def:physics:phyx-mini-0485:phyxminiproblems-problemphyxmini0485-satisfieswaterbulkmasslaw}
  \lean{PhyXMiniProblems.ProblemPhyXMini0485.SatisfiesWaterBulkMassLaw}
  \uses{def:physics:phyx-mini-0485:phyxminiproblems-problemphyxmini0485-massinkilograms, def:physics:phyx-mini-0485:phyxminiproblems-problemphyxmini0485-volumeinliters, def:physics:phyx-mini-0485:phyxminiproblems-problemphyxmini0485-massdensityinkilogramsperliter, def:physics:phyx-mini-0485:phyxminiproblems-problemphyxmini0485-horseshoewatercalorimetrysetup}
  The bulk relation m = ρV for the water in compatible unit readouts
\end{definition}
\begin{proof}
  This is the declared model definition of Satisfies Water Bulk Mass Law.
\end{proof}
\begin{definition}[Satisfies Isolated Horseshoe Water Calorimetry]
  \label{def:physics:phyx-mini-0485:phyxminiproblems-problemphyxmini0485-satisfiesisolatedhorseshoewatercalorimetry}
  \lean{PhyXMiniProblems.ProblemPhyXMini0485.SatisfiesIsolatedHorseshoeWaterCalorimetry}
  \uses{def:physics:phyx-mini-0485:phyxminiproblems-problemphyxmini0485-massinkilograms, def:physics:phyx-mini-0485:phyxminiproblems-problemphyxmini0485-specificheatinjoulesperkilogramkelvin, def:physics:phyx-mini-0485:phyxminiproblems-problemphyxmini0485-heatinjoules, def:physics:phyx-mini-0485:phyxminiproblems-problemphyxmini0485-horseshoewatercalorimetrysetup}
  The governing three-body calorimetry laws: each signed heat is mc(Tf - Ti), an isolated boundary supplies zero heat, and all four signed transfers sum to zero. These are uniform laws and do not encode the requested temperature
\end{definition}
\begin{proof}
  This is the declared model definition of Satisfies Isolated Horseshoe Water Calorimetry.
\end{proof}
\begin{definition}[Answer Choice]
  \label{def:physics:phyx-mini-0485:phyxminiproblems-problemphyxmini0485-answerchoice}
  \lean{PhyXMiniProblems.ProblemPhyXMini0485.AnswerChoice}
  Labels of the four answer choices
\end{definition}
\begin{proof}
  This is the declared model definition of Answer Choice.
\end{proof}
\begin{definition}[displayed Temperature In Degrees Celsius]
  \label{def:physics:phyx-mini-0485:phyxminiproblems-problemphyxmini0485-displayedtemperatureindegreescelsius}
  \lean{PhyXMiniProblems.ProblemPhyXMini0485.displayedTemperatureInDegreesCelsius}
  \uses{def:physics:phyx-mini-0485:phyxminiproblems-problemphyxmini0485-answerchoice}
  Celsius values printed beside the four answer labels
\end{definition}
\begin{proof}
  This is the declared model definition of displayed Temperature In Degrees Celsius.
\end{proof}
\begin{definition}[recorded Dataset Answer]
  \label{def:physics:phyx-mini-0485:phyxminiproblems-problemphyxmini0485-recordeddatasetanswer}
  \lean{PhyXMiniProblems.ProblemPhyXMini0485.recordedDatasetAnswer}
  \uses{def:physics:phyx-mini-0485:phyxminiproblems-problemphyxmini0485-answerchoice}
  The answer label recorded in the source dataset, retained as metadata
\end{definition}
\begin{proof}
  This is the declared model definition of recorded Dataset Answer.
\end{proof}
\begin{definition}[Rounds To Displayed Ten Degrees]
  \label{def:physics:phyx-mini-0485:phyxminiproblems-problemphyxmini0485-roundstodisplayedtendegrees}
  \lean{PhyXMiniProblems.ProblemPhyXMini0485.RoundsToDisplayedTenDegrees}
  \uses{def:physics:phyx-mini-0485:phyxminiproblems-problemphyxmini0485-horseshoewatercalorimetrysetup, def:physics:phyx-mini-0485:phyxminiproblems-problemphyxmini0485-initialhorseshoetemperatureindegreescelsius, def:physics:phyx-mini-0485:phyxminiproblems-problemphyxmini0485-answerchoice, def:physics:phyx-mini-0485:phyxminiproblems-problemphyxmini0485-displayedtemperatureindegreescelsius}
  The calculated temperature rounds to the displayed ten-degree value
\end{definition}
\begin{proof}
  This is the declared model definition of Rounds To Displayed Ten Degrees.
\end{proof}
\begin{definition}[Is Unique Closest Displayed Temperature]
  \label{def:physics:phyx-mini-0485:phyxminiproblems-problemphyxmini0485-isuniqueclosestdisplayedtemperature}
  \lean{PhyXMiniProblems.ProblemPhyXMini0485.IsUniqueClosestDisplayedTemperature}
  \uses{def:physics:phyx-mini-0485:phyxminiproblems-problemphyxmini0485-horseshoewatercalorimetrysetup, def:physics:phyx-mini-0485:phyxminiproblems-problemphyxmini0485-initialhorseshoetemperatureindegreescelsius, def:physics:phyx-mini-0485:phyxminiproblems-problemphyxmini0485-answerchoice, def:physics:phyx-mini-0485:phyxminiproblems-problemphyxmini0485-displayedtemperatureindegreescelsius}
  The selected displayed value is strictly closer than every alternative
\end{definition}
\begin{proof}
  This is the declared model definition of Is Unique Closest Displayed Temperature.
\end{proof}
\begin{lemma}[water mass from density and volume]
  \label{lem:physics:phyx-mini-0485:phyxminiproblems-problemphyxmini0485-water-mass-from-density-and-volume}
  \lean{PhyXMiniProblems.ProblemPhyXMini0485.water_mass_from_density_and_volume}
  \uses{def:physics:phyx-mini-0485:phyxminiproblems-problemphyxmini0485-massinkilograms, def:physics:phyx-mini-0485:phyxminiproblems-problemphyxmini0485-horseshoewatercalorimetrysetup, def:physics:phyx-mini-0485:phyxminiproblems-problemphyxmini0485-matchesproblemreadouts, def:physics:phyx-mini-0485:phyxminiproblems-problemphyxmini0485-usesreferenceironandwaterthermaldata, def:physics:phyx-mini-0485:phyxminiproblems-problemphyxmini0485-satisfieswaterbulkmasslaw}
  The density and volume readouts determine the water mass as 5/4 kg
\end{lemma}
\begin{proof}
  The conclusion follows from the governing relations and figure readouts named in the statement of water mass from density and volume.
\end{proof}
\begin{lemma}[initial horseshoe temperature exact]
  \label{lem:physics:phyx-mini-0485:phyxminiproblems-problemphyxmini0485-initial-horseshoe-temperature-exact}
  \lean{PhyXMiniProblems.ProblemPhyXMini0485.initial_horseshoe_temperature_exact}
  \uses{def:physics:phyx-mini-0485:phyxminiproblems-problemphyxmini0485-horseshoewatercalorimetrysetup, def:physics:phyx-mini-0485:phyxminiproblems-problemphyxmini0485-initialhorseshoetemperatureindegreescelsius, def:physics:phyx-mini-0485:phyxminiproblems-problemphyxmini0485-matcheshorseshoewaterscenario, def:physics:phyx-mini-0485:phyxminiproblems-problemphyxmini0485-matchesproblemreadouts, def:physics:phyx-mini-0485:phyxminiproblems-problemphyxmini0485-usesreferenceironandwaterthermaldata, def:physics:phyx-mini-0485:phyxminiproblems-problemphyxmini0485-satisfieswaterbulkmasslaw, def:physics:phyx-mini-0485:phyxminiproblems-problemphyxmini0485-satisfiesisolatedhorseshoewatercalorimetry}
  Substitution into the isolated heat balance gives 12535/72 °C, approximately 174.10 °C. The requested temperature remains on the conclusion side
\end{lemma}
\begin{proof}
  The conclusion follows from the governing relations and figure readouts named in the statement of initial horseshoe temperature exact.
\end{proof}
% --- Archon declaration topology end ---
% --- Archon physics formalization source end ---
